\title{Audible Frequency Response of Stable and Unstable Control Algorithms}
\subtitle{ECE 598 final project · University of Illinois Urbana-Champaign · Instructor: Dr.\ Arijit Banerjee}
\period{Spring 2024}
\date{2024-05}
\status{Lab developed}
\order{10}
\tags{converters, control, uiuc}

\attach{Final report}{/files/ECE 598 Final Report.docx}

\begin{document}

A teaching lab that lets you \textit{hear} a control loop go unstable. Much of
a power converter's transfer function lives inside the audible band, and the
ear reads sound as a frequency spectrum — so a boost converter wired to a
speaker turns loop stability into something you can identify by listening,
before you have read a single Bode plot.

\section{The idea}

The converter switches between two operating points once a second. Every time
it steps, the output rings according to the transfer function, and that ringing
is audible. Change the control and the sound changes with it — same circuit,
same operating point, different character.

\begin{itemize}
  \item \textbf{Open loop} --- duty cycle alternating between 20\% and 60\% at
        a fixed 10~V input, with no feedback to correct the overshoot.
  \item \textbf{Closed loop} --- reference voltage alternating between 14~V and
        18~V, with $K_p$ and $K_i$ swept from wildly unstable to critically
        damped.
  \item \textbf{Frequency sweep} --- available in either mode, walking the step
        rate from 20~Hz to 5~kHz over 5 or 10 seconds.
\end{itemize}

\section{Hardware}

A conventional boost converter on ECE 469 boards into a 50~$\Omega$ load, with
a speaker in series with a 220~$\mu$F capacitor to block the DC. Switching
signal, output current, output voltage and the mid-point sense node are all
probed, so the scope trace and the sound can be read against each other.

The closed-loop controller borrows the buck implementation and inverts the
integral term, which behaves closely enough to a true boost controller for the
purpose. The topology is incidental here — what matters is being able to dial
the stability margin freely, and that a boost output swings hard enough to
drive a speaker at a useful volume.

\section{What it sounds like}

\begin{tabular}{lll}
  Mode & Control & Sound \\
  Open loop & no feedback & a bouncing rubber ball --- impact, then reverberation \\
  Mode 0 & unstable, eventually settles & a turbine winding down, high to low \\
  Mode 1 & stable but uncontrolled & an alarm; grainy, resonant, unclear pitch \\
  Mode 2 & near the stability edge & the same pitch as Mode 1, much less sibilance \\
  Mode 3 & critically damped & near silence; a faint low thump on each step \\
\end{tabular}

Mode 1 and Mode 2 are the useful pair. Their waveforms look alike and they
measure the same loudness, so nothing on the scope separates them at a glance
--- but the resonance and grain that Mode 2 removes are immediately obvious by
ear. Mode 1 overshoots to nearly twice the intended output and still reads as a
single clean impulse, because a stable controller does not sustain the ring;
Mode 0, which is genuinely unstable, is the one that keeps going.

The sweeps make the same point on a longer time base. In Mode 2 the
peak-to-peak output grows as the step rate climbs, and the sound picks up
chirps and digital artefacts exactly where the waveform jumps. Mode 3 holds its
amplitude across the whole sweep and stays quiet throughout, apart from one
brief chirp where the sweep resets.

\section{Why it works as a lab}

The diagnostic is the point: stability margin becomes a perceptual judgement
rather than a measurement. A student who has heard Mode 1 and Mode 3 back to
back has a reference for what an underdamped loop does, and can recognise it in
a converter that has no speaker attached at all.

\end{document}
